\documentclass[12pt]{article}
\usepackage[T1]{fontenc}
\usepackage[utf8]{inputenc}
\usepackage{lmodern}
\usepackage{textcomp}
\usepackage{enumitem}
\usepackage{geometry}
\geometry{margin=1in}
\begin{document}
\begin{center}
Answer Key - Health Sciences - Graduate Level Assessment 2 \
\small (Answers are ordered exactly as in the question paper.)
\end{center}

\section*{Multiple Choice}
\begin{enumerate}[label=\arabic*.]
\item \textbf{Answer:} D
\\ \textit{Options:} A) Chylomicrons only; B) LDL only; C) VLDL and LDL; D) Chylomicrons and HDL; E) VLDL and HDL
\\ \textit{Note:} Chylomicrons and HDL are the primary lipoproteins responsible for transporting dietary cholesterol and fatty acids from the intestinal lumen, where they are absorbed by enterocytes, into the lymphatic system for circulation throughout the body.
\item \textbf{Answer:} B
\\ \textit{Options:} A) The foramen ovale; B) The foramen ovale, rotundum and spinosum; C) The foramen ovale, foramen spinosum and foramen lacerum; D) The foramen ovale and rotundum; E) The foramen ovale, rotundum, spinosum and superior orbital fissure
\\ \textit{Note:} The foramen ovale, rotundum, and spinosum are the specific foramina that transmit important nerves and vessels through the sphenoid bone, which is a critical structure in the skull that houses these openings.
\item \textbf{Answer:} B
\\ \textit{Options:} A) Afferent arteriole; B) Glomerulus; C) Renal medulla; D) Proximal convoluted tubule; E) Efferent arteriole
\\ \textit{Note:} The glomerulus is the correct answer because it specifically refers to the network of tiny blood capillaries located at the beginning of each nephron in the kidney, where the filtration of blood occurs.
\item \textbf{Answer:} D
\\ \textit{Options:} A) One gram of glycerol; B) One gram of lactose; C) One gram of fructose; D) One gram of palmitic acid; E) One gram of cellulose
\\ \textit{Note:} One gram of palmitic acid releases the most energy when completely oxidized in the body due to its long hydrocarbon chain, which contains a high number of carbon-hydrogen bonds, leading to a greater yield of ATP during cellular respiration compared to other macronutrients.
\item \textbf{Answer:} A
\\ \textit{Options:} A) Chylomicrons; B) Medium density lipoproteins; C) Micro lipoproteins; D) Intermediate density lipoproteins; E) Low density lipoproteins
\\ \textit{Note:} Chylomicrons are the lipoproteins specifically produced in the intestine after the ingestion of dietary fats, representing the exogenous lipoprotein pathway that transports lipids from the digestive system to other tissues.
\end{enumerate}

\section*{True / False}
\begin{enumerate}[label=\arabic*.]
\setcounter{enumi}{5}
\item \textbf{Answer:} False
\\ \textit{Options:} A) True; B) False
\\ \textit{Note:} Vitamin E is primarily an antioxidant and does not function as a coenzyme in fatty acid synthesis, where coenzymes like NADPH play a crucial role in reduction reactions.
\item \textbf{Answer:} False
\\ \textit{Options:} A) True; B) False
\\ \textit{Note:} The stimulation of appetite is not the most important gastrointestinal function for survival immediately after a meal because digestion and nutrient absorption are crucial for providing energy and sustaining bodily functions, rather than merely stimulating further appetite.
\item \textbf{Answer:} True
\\ \textit{Options:} A) True; B) False
\\ \textit{Note:} Most pandemics have originated from influenza viruses carried by wild birds due to the ability of these viruses to mutate and jump from avian hosts to humans, leading to novel strains that can cause widespread illness.
\item \textbf{Answer:} True
\\ \textit{Options:} A) True; B) False
\\ \textit{Note:} The complete resynthesis of phosphocreatine after very high intensity exercise takes about 4 minutes due to the rapid recovery process that occurs through aerobic metabolism and the replenishment of creatine phosphate stores in muscle tissue.
\item \textbf{Answer:} True
\\ \textit{Options:} A) True; B) False
\\ \textit{Note:} The statement is true because medical guidelines emphasize that clinicians should minimize patient discomfort and complications by seeking assistance after one unsuccessful cannulation attempt, thereby ensuring patient safety and adherence to best practices in clinical care.
\end{enumerate}

\section*{Long Form Response}
\begin{enumerate}[label=\arabic*.]
\setcounter{enumi}{10}
\item \textbf{Answer:} Excessive consumption of Vitamin C, while generally considered safe due to its water- soluble nature, can lead to gastrointestinal disturbances, such as diarrhea and abdominal cramps, particularly at doses exceeding 2,000 mg per day. In contrast, other food supplements, like fat-soluble vitamins (e.g., A, D, E, K), pose a higher risk of toxicity because they accumulate in body tissues. However, Vitamin C does offer significant benefits, such as enhanced immune function and antioxidant properties, which can contribute to overall health when consumed in appropriate amounts. The key distinction lies in the balance between benefits and risks; while moderate intake of Vitamin C can be beneficial, excessive amounts may lead to adverse effects, necessitating careful consideration of dietary sources and supplementation strategies. In comparison, other supplements may require more stringent monitoring to avoid toxicity, emphasizing the importance of individualized approaches to nutrient supplementation.
\\ \textit{Note:} The answer correctly highlights that while excessive Vitamin C intake can cause gastrointestinal issues, it poses a lower risk of toxicity compared to fat-soluble vitamins, while also acknowledging the significant health benefits associated with Vitamin C, thus providing a balanced view of both its risks and advantages.
\item \textbf{Answer:} Research indicates that the prevalence of respiratory diseases tends to be lower among vegetarians compared to non-vegetarians, a trend potentially influenced by several lifestyle and dietary factors. Vegetarians often consume a diet rich in fruits, vegetables, and whole grains, which are high in antioxidants and anti-inflammatory compounds that can enhance lung health. Additionally, vegetarians are less likely to smoke and may engage in more health-promoting behaviors, such as regular physical activity, contributing to better respiratory function. Conversely, non-vegetarians may have higher exposure to certain dietary triggers, such as processed meats, which have been linked to increased inflammation and respiratory issues. Therefore, while dietary choices play a crucial role, broader lifestyle factors also significantly impact the prevalence of respiratory diseases in these populations.
\\ \textit{Note:} correct because it highlights that a vegetarian diet, which is typically high in antioxidants and anti-inflammatory foods, along with healthier lifestyle choices such as lower smoking rates, contributes to a reduced prevalence of respiratory diseases compared to non-vegetarians.
\end{enumerate}

\vfill
\noindent\textit{End of Answer Key}
\end{document}